%=============================================================================!
% 14.0  CHAPTER OPENING                                                       !
%=============================================================================!
A graphite electrode swells as it charges: the carbon sheets separate
from 3.34 to \SI{3.70}{\angstrom} when lithium fills the spaces between
them\cite{Hazrati2014}. Holding a simulation at the volume of the bare
graphite would compress the lithiated solid across its sheets. To study
its response at a prescribed pressure, we must allow the cell's volume
and shape to change.

We first obtain pressure from the force of atoms on a wall, then use the
virial to calculate it in a periodic cell without walls. The pressure
tensor distinguishes the forces on different faces, allowing us to
describe a layered solid whose response differs within and across its
sheets. We then develop barostats and test both the mean pressure and the
volume fluctuations they produce, following the approach to thermostats
in Chapter~13. Their effects on motion and their characteristic faults
prepare the question of Chapter~15: how long must a run be before its
averages can be trusted?
